\section{Relazioni tra insiemi}
Ancora una volta mi vedo costretto a rimarcare concetti basilari, al solo
scopo di avere un testo pulito e chiaro. \\
Procederò con lo spiegare cosa è una \textbf{funzione}.
\\
\\
Abbiamo parlato ampiamente di insiemi, cosa sono, come si comportano quando
soggetti a delle operazioni e tanto altro. Adesso però parleremo delle 
\textbf{relazioni tra insiemi}. \\
Prendiamo l'insieme \(Cesto = \{Mela, Arancia, Banana\}\), ovvero la frutta
all'interno di un cesto, e l'insieme 
\( Colori = \{Rosso, Arancione, Giallo, Verde\}\),
una semplice raccolta di colori. Definiamo allora una \textbf{funzione} \(f\)
che "mette in relazione" i due insieme, in questo caso, defineremo \(f\) come
\[ f(\text{frutta}) = \text{colore di frutta} \]
Usando i due insiemi appena definiti vedremo che la nostra funzione si comporta
nel seguente modo:
\begin{align*}
  f(\text{Mela}) &= \text{Rosso} \\
  f(\text{Arancia}) &= \text{Arancione} \\
  f(\text{Banana}) &= \text{Giallo}
\end{align*} 
La nostra funzione \(f\) ha \emph{messo in relazione} \underline{ciascun}
elemento dell'insieme \(Cesto\) con un altro dell'insieme \(Colori\); notare
come non tutti gli elementi di \(Colori\) sono "collegati" a elementi di 
\(Cesto\), la cosa importante è che \underline{ogni} elemento del primo 
insieme sia collegato ad almeno un elemento del secondo. \\
Abbiamo appena definito un \textbf{morfismo}.
\\
\\
In modo generale diremo quindi che un \textbf{morfismo} \(f\) è un processo che
trasforma una struttura di partenza \(A\) in una diversa strttura \(B\); nulla
toglie al processo di "trasformare" una strttura in sè stessa. \\
Di seguito vederemo vari tipi di morfismi.

\subsection{Omomorfismi}
Come detto prima un morfismo è una trasformazione riguardanti strtture;
ma questo termine è molto generico. \\
Indicheremo quindi col termine \textbf{omomorfismo} quel tipo di morfismo 
riguardante \textbf{strtture algebriche}. La definizione è praticamente la
stessa di quella definita prima, quindi passerò direttamente a fare un esempio.
\\
\\
Sia \(f: \mathbb{N} \to \mathbb{Z} \), ovvero una funzione che associa 
ogni elemento dell'insieme dei numeri naturali con un elemento 
dei numeri interi, e sia \(f = x \cdot -1\), ovvero una funzione che ad ogni
elemento di \(\mathbb{N}\) moltiplica \(-1\). \\
